\section{Open Questions, Ambiguities, and Practical Roadmap}

\subsection{Ambiguity register}

Several concepts are clearly important in the materials but remain too fuzzy to encode directly.

\begin{itemize}[leftmargin=1.4em]
  \item \term{Clean versus dirty setup quality.} Traders recognize it quickly, but the boundary is
        rarely formalized.
  \item \term{Real participant versus fake liquidity.} The definition is behavioral, not static.
  \item \term{Robot identification.} Some cases are obvious to a trained eye, but pattern boundaries
        are unclear.
  \item \term{Smart Money.} Useful framing, weakly operational on its own.
  \item \term{Volatility harvesting.} Mentioned conceptually, weakly specified in local evidence.
  \item \term{Market regime transfer.} The same visible pattern can mean opposite things depending on
        whether the day is trend, rotation, squeeze, or event-driven chaos.
\end{itemize}

\subsection{Possible contradictions or source mismatches}

The knowledge base mixes educational books and practice-first transcript lessons. The books are
cleaner and more explanatory. The transcripts are more operational but also more heuristic. This
creates expected tension:

\begin{itemize}[leftmargin=1.4em]
  \item Books can make setups appear more stable and teachable than they feel in live examples.
  \item Transcripts can sound highly discretionary and sometimes downplay the need for formal
        repeatability.
  \item Risk books emphasize conservative structure, while live examples sometimes rely on very fast
        behavioral exits that are hard to reproduce without high-quality execution.
\end{itemize}

These are not fatal contradictions. They simply mean the framework should record both explicit
rules and a separate ambiguity register.

\subsection{Questions that require empirical testing}

\begin{enumerate}[leftmargin=1.4em]
  \item Which density metrics are most predictive once normalized by instrument and regime?
  \item How often do visible density bounces survive once trades / prints are incorporated?
  \item Can breakout quality be separated reliably using only short-horizon order-flow features?
  \item How stable are leader-follower relations within one day and across days?
  \item Can same-day calibration materially improve rule-based detectors without overfitting?
  \item Which setups remain meaningful in crypto versus Russian markets, and which are venue-specific?
\end{enumerate}

\subsection{Practical roadmap for the repository}

\subsubsection{Phase 1: formalize and log}

Complete the knowledge formalization, keep the viewer tools organized, and use the new Bybit and
T-Invest logging layer to start collecting research-grade data. During this phase, the win
condition is not profit. The win condition is better evidence.

\subsubsection{Phase 2: test simple claims}

Use replay and live capture to test the most basic relationships:

\begin{itemize}[leftmargin=1.4em]
  \item density bounce frequency,
  \item density pull plus breakout,
  \item failed continuation after breakout,
  \item reaction around visible absorption,
  \item early-session pattern quality.
\end{itemize}

\subsubsection{Phase 3: encode rule-based setup detectors}

Build detectors that are explicit, configurable, and reviewable. They should emit not only
signals but also \emph{reason bundles}: the components that caused the setup to be flagged.

\subsubsection{Phase 4: add replay, evaluation, and analytics}

The platform should compare:

\begin{itemize}[leftmargin=1.4em]
  \item historical replay assumptions,
  \item live observed behavior,
  \item signal quality by instrument,
  \item signal quality by time of day,
  \item signal quality by day regime.
\end{itemize}

\subsubsection{Phase 5: only then consider adaptive filters}

Once the repository has richer live data and rule-based events, it becomes reasonable to
consider ML as a \emph{filter} rather than as an oracle. The most promising uses are:

\begin{itemize}[leftmargin=1.4em]
  \item taking or skipping borderline setups,
  \item classifying breakout versus fakeout,
  \item ranking setup cleanliness,
  \item adapting to current-day context.
\end{itemize}

\subsection{Final working conclusion}

The central message of the knowledge base is not ``watch big orders.'' The central message is
that short-term market behavior is a negotiation between visible liquidity, hidden liquidity,
aggressive flow, nearby structural context, and changing day conditions. A useful research
platform must therefore preserve that negotiation in data form. Once the negotiation is
preserved, the repository can move gradually from discretionary stories toward explicit event
patterns, measured features, and finally robust live-tested logic.
